\section{Test Case 16: dynamic Maxwell evaporation on the GPU}

Test Case 16 exercises Maxwell rheology, Yarin evaporation, nozzle insertion,
collector removal, and capacity growth in one serial three-dimensional run.
It starts with 100 beads and deliberately retains the small initial capacity
so that the 1,000-step run invokes the reallocation path.

The stored input selects RK4; changing only \texttt{integrator} to 1 or 2
selects the Euler or RK2 validation.  With target
\texttt{nvfortran-openacc}, Euler executes one Maxwell force and stress
evaluation, RK2 executes two, and RK4 executes four.  Their intermediate and
final updates, direct Coulomb summation, evaporation, charge smoothing and
restoration, nozzle geometry, and final statistics reductions execute on the
GPU.  No state, force, stress, or derivative array is transferred between
stages.  Ordinary timesteps return only topology decision scalars.  Bead
records are moved for actual topology or output events, while the complete
active state is synchronized for capacity reallocation and the final
checkpoint.

On an NVIDIA A30 the accepted 1,000-step topology for Euler, RK2, and RK4
contains 111 additions, 122 removals, two reallocations, and 89 final active
beads, equal to the CPU totals.  For Euler, RK2, and RK4, three-step pre-event
\texttt{statout.dat} comparisons have zero difference with relative tolerance
$10^{-12}$ and absolute tolerance $10^{-13}$; their written XYZ geometries are
also byte-identical.  Floating-point ordering moves the first insertion from
step 4 on the CPU to step 5 on the GPU.  A pointwise CPU/GPU trajectory after
this threshold is not used as the acceptance criterion: the target-centric
direct Coulomb kernel has a different summation order, and the physical
bending instability amplifies the resulting roundoff-level perturbation.

The rheology-independent development target
\texttt{nvfortran-openacc-force-oracle} evaluates each complete trusted
Maxwell force stage on the host and uploads its derivatives while state
updates remain on the GPU.  The narrower
\texttt{nvfortran-openacc-coulomb-oracle} target evaluates only the direct
Coulomb sum on the host. Both reproduce the accepted 1,000-step topology for
Euler, RK2, and RK4, but their deliberate transfers make them unsuitable for
production or performance runs. The input is stored in
\texttt{examples/input-16/input.dat}.

The local script
\texttt{tests/performance/dynamic/validate\_evaporation.sh} builds independent
NVFORTRAN CPU and standard OpenACC executables and repeats the complete paired
Euler, RK2, and RK4 validation for Test Cases 16 and 17.  The standard GPU
build used by the script does not enable either diagnostic oracle macro.
